% s-220-p220-m003; original PDF 220, 221
\Aphorism{s-220-p220-m003}{314}
Um eine Vergangenheit abzuschütteln, kanonisirt man seine Gegenwart — und beschimpft damit vielleicht seine eigene Zukunft! Man ahnt nie, wie sehr sich die Entwicklung eines modernen Geistes der Kreisform, mindestens der Spiralform nähern kann; man verlässt eine Cultur durch das Hauptportal, um vielleicht durch ein Hinterpförtchen wieder hineinzuschlüpfen. Es giebt eine Periodik der Zustände, ein Hin und Her der Metamorphosen; der Enkel schlägt mehr dem Grossvater nach als dem Vater, die dritte Weltanschauung, conträr zur zweiten, neigt wieder zur ersten, der Schmetterling, mit Überspringung der Larve, ähnelt der Raupe. Indem man sich selbst vergewaltigt, sich in jedem Augenblick stabilisirt und auf eine Formel, eine gegen das Vergangene polemische Formel bringt, begeht man eine Ungerechtigkeit, gegen die vielleicht schon die nächste Zukunft mit einer neuen Ungerechtigkeit protestiren muss; jede dieser Pendelschwingungen zieht mit Nothwendigkeit die folgenden nach sich. — Nietzsches „Buch für freie Geister“ ist eine grundsätzliche Apotheose auf das, was man klassischen Stil zu nennen liebt: seine Anbetung vor den Formen, vor Mitte und Mass, vor der kühlen, schlichten, reizlosen Darstellungsweise einiger Griechen und Franzosen, seine Abneigung gegen das „Genie“, gegen jeden unüberwundenen Rest von Kraft, Gefühl und Orgiasmus geht ins Asketische. Man weiss, warum; der mittlere Nietzsche fällt vom jugendlichen ab. Aber wenn dieses menschlich-allzumenschliche Buch an sich Recht hätte, so behielte auch der spätere Nietzsche Unrecht, und vielleicht gründlicher als der frühere. Die Merkmale des Barockstiles, die hier der Kunst zweiten und dritten Ranges vorgerechnet werden, das Gewaltsame, Explosive, leidenschaftlich Bunte, die „ Dämmerungs- Feuersbrunst- und Verklärungslichter auf stark-gebildeten Formen“, das Hautrelief in Abstufungen und Gegensätzen — dies alles lässt sich ebensogut am „Jenseits von Gut und Böse“ studiren wie an der „ Geburt der Tragödie“ und Wagnerischer Musik; der dithyrambische Sturmhauch der Zarathustra-Dichtung ist nicht so sehr unterschieden von der ekstatischen Inbrunst in „Tristan und Isolde“, wie beide von der positivistischen Kälte und Sachlichkeit des „freien Geistes“ unterschieden sind. Man bemüht sich umsonst, die immer erneute Parteinahme gegen sich selbst, die schmerzliche Stufenfolge von Selbstentzweiung und Selbstüberwindung aus Nietzsches Entwicklung fortzudeuten; der bei modernen Seelen ohnehin schon so hoch hinauf getriebene „innere Widerstand“ ist bei ihm ein Maximum — daher die Widersprüche vor- und rückwärts, die beständige Überschreitung der Gleichgewichtslage nach der einen oder anderen Seite, die hin und her oscillirenden Entladungen einer ungeheuer hochgespannten Polarität.
\par

